\renewcommand{\currentchapter}{7}
\title{Contoh Soal UTBK}
\subtitle{Bab 7 \textbar\ Tujuh set representatif, pembahasan, dan strategi cepat}
\date{SMAN 1 Muntilan bersama StudentPro \textbar\ Tahun Ajaran 2026/2027}

\newcommand{\examplemeta}[3]{%
  {\color{StudentBlue}\bfseries\fontsize{10.4}{12.8}\selectfont #1}%
  {\color{StudentGray}\fontsize{10.2}{12.8}\selectfont\enspace\textbar\enspace #2\enspace\textbar\enspace #3}\par\vspace{0.10cm}%
}
\newcommand{\mistake}[1]{%
  {\color{StudentRed}\bfseries Kesalahan umum:}\enspace #1%
}
\newcommand{\quick}[1]{%
  {\color{StudentBlue}\bfseries Strategi cepat:}\enspace #1%
}

\begingroup
\setbeamercolor{background canvas}{bg=StudentNavy}
\begin{frame}[plain]
  \vspace{0.65cm}
  \kicker{SMAN 1 Muntilan \textbar\ StudentPro}
  \vspace{0.08cm}
  {\usebeamerfont{title}\color{StudentWhite}\inserttitle\par}
  \vspace{0.38cm}
  {\usebeamerfont{subtitle}\color{StudentSoft}\insertsubtitle\par}
  \vfill
  {\color{StudentGold}\rule{2.3cm}{2.5pt}}\par
  \vspace{0.20cm}
  {\usebeamerfont{author}\color{StudentWhite}\insertauthor}\par
  {\usebeamerfont{date}\color{StudentSoft}\insertdate}
  \vspace{0.42cm}
\end{frame}
\endgroup

\begin{frame}{Bab 7 menunjukkan cara kemampuan diuji}
  \kicker{Cara Membaca Contoh}
  \begin{columns}[T,onlytextwidth]
    \begin{column}{0.30\textwidth}
      \begin{block}{Stimulus}
        Bacaan, tabel, aturan, atau kasus yang harus dipahami terlebih dahulu.
      \end{block}
    \end{column}
    \begin{column}{0.30\textwidth}
      \begin{block}{Keputusan}
        Siswa memilih jawaban berdasarkan bukti, hubungan, atau model.
      \end{block}
    \end{column}
    \begin{column}{0.30\textwidth}
      \begin{block}{Refleksi}
        Jawaban dibedah melalui alasan, kesalahan umum, dan strategi cepat.
      \end{block}
    \end{column}
  \end{columns}
  \vspace{0.25cm}
  \begin{beamercolorbox}[rounded=true,sep=0.18cm,wd=\textwidth]{block body}
    \centering\bfseries Seluruh contoh disusun orisinal oleh StudentPro; bukan soal resmi atau soal bocoran SNPMB.
  \end{beamercolorbox}
\end{frame}

\begin{frame}[shrink=8]{Tujuh set menampilkan variasi tuntutan berpikir}
  \kicker{Peta Contoh}
  {\fontsize{9.7}{12.0}\selectfont
  \renewcommand{\arraystretch}{0.96}
  \begin{tabularx}{\textwidth}{>{\bfseries\color{StudentBlue}}p{0.25\textwidth} >{\RaggedRight\arraybackslash}p{0.23\textwidth} >{\RaggedRight\arraybackslash}X}
    \toprule
    \textbf{Subtes} & \textbf{Stimulus/bentuk} & \textbf{Kemampuan yang ditunjukkan} \\
    \midrule
    Penalaran Umum & Aturan dan kasus & Menentukan simpulan yang wajib benar \\
    PPU & Paragraf singkat & Memaknai kata menurut konteks \\
    PBM & Kalimat bermasalah & Memilih perbaikan paling efektif \\
    Pengetahuan Kuantitatif & Komparasi $P$--$Q$ & Menggunakan identitas secara efisien \\
    LBI & Tabel dan dua bacaan & Menilai bukti serta merefleksikan teks \\
    LBE & Academic passage & Menemukan klaim yang didukung bacaan \\
    Penalaran Matematika & Kecukupan data & Merumuskan model dari dua informasi \\
    \bottomrule
  \end{tabularx}}
  \vspace{0.10cm}
  {\color{StudentGray}\fontsize{9.7}{12}\selectfont Variasi bentuk merupakan rancangan latihan StudentPro, bukan kepastian format resmi UTBK 2027.}
\end{frame}

\begin{frame}[shrink=10]{Contoh 1 -- Penalaran Umum}
  \kicker{Soal \textbar\ Deduktif}
  \examplemeta{Indikator: simpulan valid}{Kesulitan: sedang}{Target: 75 detik}
  \begin{block}{Aturan kegiatan}
    Semua peserta Klinik Data mengikuti diagnostik. Hanya peserta yang tuntas diagnostik yang boleh memasuki simulasi. Setiap peserta simulasi wajib mengikuti evaluasi. Raka mengikuti Klinik Data, tetapi belum tuntas diagnostik.
  \end{block}
  {\bfseries Simpulan yang pasti benar adalah~\ldots}
  \vspace{0.10cm}
  {\fontsize{10.1}{12.5}\selectfont
  \begin{columns}[T,onlytextwidth]
    \begin{column}{0.48\textwidth}
      A. Raka tidak mengikuti diagnostik.\par\vspace{0.08cm}
      B. Raka belum boleh memasuki simulasi.\par\vspace{0.08cm}
      C. Raka pasti tidak mengikuti evaluasi.
    \end{column}
    \begin{column}{0.48\textwidth}
      D. Semua peserta diagnostik mengikuti Klinik Data.\par\vspace{0.08cm}
      E. Semua peserta evaluasi telah tuntas diagnostik.
    \end{column}
  \end{columns}}
\end{frame}

\begin{frame}{Bedah Contoh 1 -- Penalaran Umum}
  \kicker{Jawaban dan Proses Berpikir}
  \answer{B. Raka belum boleh memasuki simulasi.}
  \vspace{0.22cm}
  \begin{block}{Mengapa?}
    Klinik Data memastikan Raka mengikuti diagnostik. Aturan ``hanya yang tuntas'' menjadikan ketuntasan syarat wajib untuk masuk simulasi. Karena Raka belum tuntas, ia belum boleh masuk.
  \end{block}
  {\fontsize{10.6}{13.2}\selectfont
  \mistake{membalik hubungan ``semua peserta Klinik mengikuti diagnostik'' menjadi ``semua peserta diagnostik mengikuti Klinik''.}\par\vspace{0.16cm}
  \quick{ubah setiap aturan menjadi panah satu arah, lalu ikuti hanya jalur yang berkaitan dengan Raka.}}
\end{frame}

\begin{frame}[shrink=6]{Contoh 2 -- Pengetahuan dan Pemahaman Umum}
  \kicker{Soal \textbar\ Makna Kontekstual}
  \examplemeta{Indikator: makna kata}{Kesulitan: sedang}{Target: 60 detik}
  \begin{block}{Stimulus}
    Ketika hasil satu tryout menurun, tim tidak serta-merta mengubah target. Data itu diperlakukan sebagai \textbf{\emph{petunjuk}}, bukan \textbf{\emph{vonis}}. Tim menelaah jenis kesalahan, waktu pengerjaan, dan kondisi siswa sebelum menentukan tindak lanjut.
  \end{block}
  {\bfseries Makna kata \emph{vonis} dalam paragraf tersebut adalah~\ldots}
  \vspace{0.10cm}
  {\fontsize{10.1}{12.5}\selectfont
  \begin{columns}[T,onlytextwidth]
    \begin{column}{0.48\textwidth}
      A. keputusan akhir yang dianggap menentukan\par\vspace{0.08cm}
      B. bukti awal yang perlu diperiksa\par\vspace{0.08cm}
      C. catatan teknis selama latihan
    \end{column}
    \begin{column}{0.48\textwidth}
      D. perbandingan hasil antarsiswa\par\vspace{0.08cm}
      E. rencana belajar jangka pendek
    \end{column}
  \end{columns}}
\end{frame}

\begin{frame}{Bedah Contoh 2 -- PPU}
  \kicker{Jawaban dan Proses Berpikir}
  \answer{A. keputusan akhir yang dianggap menentukan.}
  \vspace{0.22cm}
  \begin{block}{Mengapa?}
    Kata \emph{vonis} dikontraskan dengan \emph{petunjuk}. Paragraf menolak menjadikan satu nilai sebagai keputusan final dan memilih memeriksa data pendukung terlebih dahulu.
  \end{block}
  {\fontsize{10.6}{13.2}\selectfont
  \mistake{memilih arti kamus yang paling dikenal tanpa menguji hubungan antarkalimat.}\par\vspace{0.16cm}
  \quick{cari penanda kontras seperti ``bukan'', lalu ganti kata sasaran dengan setiap opsi dan uji kecocokannya.}}
\end{frame}

\begin{frame}[shrink=8]{Contoh 3 -- Pemahaman Bacaan dan Menulis}
  \kicker{Soal \textbar\ Kalimat Efektif}
  \examplemeta{Indikator: penyuntingan}{Kesulitan: sedang}{Target: 60 detik}
  \begin{block}{Kalimat asal}
    Bagi siswa yang belum stabil akurasinya maka perlu diberikan latihan bertahap agar supaya mereka dapat memperbaiki ketelitian.
  \end{block}
  {\bfseries Perbaikan paling efektif adalah~\ldots}\par
  \vspace{0.08cm}
  {\fontsize{9.65}{11.9}\selectfont
  A. Bagi siswa yang akurasinya belum stabil perlu diberikan latihan bertahap agar lebih teliti.\par\vspace{0.06cm}
  B. Siswa yang akurasinya belum stabil perlu diberi latihan bertahap agar lebih teliti.\par\vspace{0.06cm}
  C. Siswa dengan akurasi belum stabil maka perlu diberikan latihan bertahap agar supaya teliti.\par\vspace{0.06cm}
  D. Latihan bertahap perlu diberikan bagi siswa yang belum stabil akurasinya supaya agar teliti.\par\vspace{0.06cm}
  E. Untuk siswa yang belum stabil akurasinya adalah perlu latihan bertahap untuk memperbaiki ketelitian.}
\end{frame}

\begin{frame}{Bedah Contoh 3 -- PBM}
  \kicker{Jawaban dan Proses Berpikir}
  \answer{B. Siswa yang akurasinya belum stabil perlu diberi latihan bertahap agar lebih teliti.}
  \vspace{0.22cm}
  \begin{block}{Mengapa?}
    Opsi B memiliki subjek dan predikat yang jelas, mempertahankan makna, serta menghilangkan pasangan mubazir \emph{bagi--maka} dan \emph{agar--supaya}.
  \end{block}
  {\fontsize{10.6}{13.2}\selectfont
  \mistake{hanya memperbaiki satu kata, tetapi membiarkan struktur tanpa subjek atau pengulangan makna.}\par\vspace{0.16cm}
  \quick{cek empat hal: subjek--predikat, ketepatan makna, kehematan, dan kelancaran saat dibaca.}}
\end{frame}

\begin{frame}[shrink=4]{Contoh 4 -- Pengetahuan Kuantitatif}
  \kicker{Soal \textbar\ Komparasi Kuantitatif}
  \examplemeta{Indikator: bentuk aljabar}{Kesulitan: sedang}{Target: 75 detik}
  \begin{block}{Informasi}
    Bilangan positif $x$ dan $y$ memenuhi $x+y=10$ dan $xy=21$.
  \end{block}
  \begin{columns}[T,onlytextwidth]
    \begin{column}{0.46\textwidth}
      \begin{block}{Kuantitas $P$}
        \centering $P=x^2+y^2$
      \end{block}
    \end{column}
    \begin{column}{0.46\textwidth}
      \begin{block}{Kuantitas $Q$}
        \centering $Q=6(x+y)-2$
      \end{block}
    \end{column}
  \end{columns}
  {\bfseries Hubungan yang benar adalah~\ldots}\par\vspace{0.09cm}
  {\fontsize{10.3}{12.8}\selectfont
  A. $P>Q$\quad B. $P<Q$\quad C. $P=Q$\quad D. $P=2Q$\quad E. tidak dapat ditentukan}
\end{frame}

\begin{frame}{Bedah Contoh 4 -- Pengetahuan Kuantitatif}
  \kicker{Jawaban dan Proses Berpikir}
  \answer{C. $P=Q$.}
  \vspace{0.18cm}
  \begin{columns}[T,onlytextwidth]
    \begin{column}{0.47\textwidth}
      \begin{block}{Hitung $P$}
        $x^2+y^2=(x+y)^2-2xy=10^2-2(21)=58$.
      \end{block}
    \end{column}
    \begin{column}{0.47\textwidth}
      \begin{block}{Hitung $Q$}
        $Q=6(x+y)-2=6(10)-2=58$.
      \end{block}
    \end{column}
  \end{columns}
  {\fontsize{10.6}{13.2}\selectfont
  \mistake{mencari nilai $x$ dan $y$ satu per satu, padahal yang dibutuhkan dapat diperoleh langsung dari jumlah dan hasil kali.}\par\vspace{0.14cm}
  \quick{sebelum menghitung, cocokkan bentuk yang ditanya dengan identitas yang memanfaatkan data secara langsung.}}
\end{frame}

\begin{frame}{LBI dibaca melalui tiga konteks}
  \kicker{Contoh 5 -- Literasi dalam Bahasa Indonesia}
  \begin{columns}[T,onlytextwidth]
    \begin{column}{0.30\textwidth}
      \begin{block}{Saintek}
        Membaca tabel, menilai desain pengamatan, dan membatasi klaim sebab-akibat.
      \end{block}
    \end{column}
    \begin{column}{0.30\textwidth}
      \begin{block}{Sosial-humaniora}
        Menilai argumen kebijakan, bukti sosial, dan penjelasan alternatif.
      \end{block}
    \end{column}
    \begin{column}{0.30\textwidth}
      \begin{block}{Personal}
        Mengapresiasi bahasa, konflik tokoh, sudut pandang, dan makna tersirat.
      \end{block}
    \end{column}
  \end{columns}
  \vspace{0.25cm}
  \claim{Ketiganya menguji proses literasi yang sama melalui jenis bacaan yang berbeda.}
\end{frame}

\begin{frame}[shrink=10]{Contoh 5A -- LBI konteks saintek}
  \kicker{Soal \textbar\ Menilai Bukti}
  \examplemeta{Indikator: evaluasi klaim}{Kesulitan: tinggi}{Target: 100 detik}
  {\fontsize{9.7}{11.8}\selectfont
  \begin{tabularx}{\textwidth}{>{\bfseries}X c c c}
    \toprule
    \textbf{Kelompok belajar} & \textbf{Skor awal} & \textbf{Skor akhir} & \textbf{Kehadiran} \\
    \midrule
    Metode A ($n=40$) & 62 & 72 & 88\% \\
    Metode B ($n=40$) & 68 & 74 & 89\% \\
    \bottomrule
  \end{tabularx}}
  \vspace{0.10cm}
  {\bfseries Penilaian paling tepat terhadap klaim ``Metode A pasti lebih efektif'' adalah~\ldots}\par
  \vspace{0.07cm}
  {\fontsize{8.9}{10.8}\selectfont
  A. Benar, karena skor akhir A lebih tinggi daripada skor awalnya.\par\vspace{0.035cm}
  B. Benar, karena kenaikan A lebih besar dan jumlah peserta sama.\par\vspace{0.035cm}
  C. Belum pasti; kenaikan A lebih besar, tetapi kesetaraan kelompok dan faktor lain belum diketahui.\par\vspace{0.035cm}
  D. Salah, karena skor akhir B tetap lebih tinggi daripada A.\par\vspace{0.035cm}
  E. Salah, karena kehadiran kedua kelompok hampir sama.}
\end{frame}

\begin{frame}{Bedah Contoh 5A -- LBI saintek}
  \kicker{Jawaban dan Proses Berpikir}
  \answer{C. Klaim kausal belum dapat dipastikan.}
  \vspace{0.20cm}
  \begin{block}{Mengapa?}
    Data menunjukkan kenaikan rata-rata A sebesar 10 poin dan B sebesar 6 poin. Namun, tabel tidak menjelaskan cara pembagian kelompok, variasi skor, pengajar, atau faktor lain yang dapat memengaruhi hasil.
  \end{block}
  {\fontsize{10.5}{13.1}\selectfont
  \mistake{menganggap selisih rata-rata otomatis membuktikan metode sebagai penyebab.}\par\vspace{0.15cm}
  \quick{pisahkan apa yang benar-benar tampak pada data dari klaim sebab-akibat yang memerlukan desain pembanding.}}
\end{frame}

\begin{frame}[shrink=10]{Contoh 5B -- LBI konteks sosial-humaniora}
  \kicker{Soal \textbar\ Evaluasi Argumen}
  \examplemeta{Indikator: bukti penguat}{Kesulitan: tinggi}{Target: 90 detik}
  \begin{block}{Stimulus}
    Setelah perpustakaan kota memperpanjang jam layanan malam, jumlah kunjungan bulanan naik 18\%. Pengelola menyimpulkan bahwa perpanjangan jam layanan menyebabkan kenaikan tersebut. Pada bulan yang sama, halte angkutan umum baru mulai beroperasi tepat di depan perpustakaan.
  \end{block}
  {\bfseries Data tambahan yang paling membantu menguji simpulan pengelola adalah~\ldots}\par
  \vspace{0.07cm}
  {\fontsize{9.3}{11.4}\selectfont
  A. jumlah koleksi buku yang dimiliki perpustakaan\par\vspace{0.04cm}
  B. warna dan tata letak ruang baca malam\par\vspace{0.04cm}
  C. kunjungan pada jam tambahan serta asal moda transportasi pengunjung\par\vspace{0.04cm}
  D. jumlah pegawai yang menyukai jadwal malam\par\vspace{0.04cm}
  E. pendapat pengelola tentang budaya membaca}
\end{frame}

\begin{frame}{Bedah Contoh 5B -- LBI sosial-humaniora}
  \kicker{Jawaban dan Proses Berpikir}
  \answer{C. Kunjungan pada jam tambahan dan moda transportasi.}
  \vspace{0.20cm}
  \begin{block}{Mengapa?}
    Pilihan C membantu membedakan dua penjelasan yang muncul bersamaan: tambahan waktu layanan dan kemudahan akses melalui halte baru.
  \end{block}
  {\fontsize{10.6}{13.2}\selectfont
  \mistake{memilih data yang berhubungan dengan perpustakaan, tetapi tidak dapat menguji penyebab kenaikan kunjungan.}\par\vspace{0.16cm}
  \quick{tentukan klaim, cari penjelasan alternatif, lalu pilih data yang dapat membandingkan keduanya.}}
\end{frame}

\begin{frame}[shrink=10]{Contoh 5C -- LBI konteks personal}
  \kicker{Soal \textbar\ Makna Tersirat}
  \examplemeta{Indikator: refleksi tokoh}{Kesulitan: sedang}{Target: 75 detik}
  \begin{block}{Kutipan fiksi orisinal}
    Nara berdiri di depan papan pengumuman. Namanya ada di urutan terakhir peserta yang lolos. Ia tersenyum, lalu melipat surat hasil tryout yang sejak pagi diremasnya. ``Ternyata pintunya tidak mengecil,'' gumamnya, ``aku saja yang terlalu lama menunduk.''
  \end{block}
  {\bfseries Makna perubahan sikap Nara paling tepat adalah~\ldots}\par
  \vspace{0.07cm}
  {\fontsize{9.6}{11.8}\selectfont
  A. Ia menilai keberhasilan hanya ditentukan keberuntungan.\par\vspace{0.04cm}
  B. Ia menyadari keraguan diri membatasi cara melihat peluang.\par\vspace{0.04cm}
  C. Ia kecewa karena berada pada urutan terakhir.\par\vspace{0.04cm}
  D. Ia ingin menyembunyikan hasil tryout dari orang lain.\par\vspace{0.04cm}
  E. Ia menganggap proses belajar tidak lagi diperlukan.}
\end{frame}

\begin{frame}{Bedah Contoh 5C -- LBI personal}
  \kicker{Jawaban dan Proses Berpikir}
  \answer{B. Keraguan diri membatasi cara Nara melihat peluang.}
  \vspace{0.20cm}
  \begin{block}{Mengapa?}
    ``Pintu'' melambangkan peluang. Pernyataan bahwa dirinya terlalu lama menunduk menunjukkan perubahan dari rasa ragu menuju kesadaran dan keberanian melihat kesempatan.
  \end{block}
  {\fontsize{10.6}{13.2}\selectfont
  \mistake{membaca ungkapan metaforis secara harfiah atau hanya berfokus pada urutan terakhir.}\par\vspace{0.16cm}
  \quick{hubungkan simbol dengan tindakan tokoh sebelum dan sesudah titik perubahan.}}
\end{frame}

\begin{frame}[shrink=10]{Contoh 6 -- Literasi dalam Bahasa Inggris}
  \kicker{Soal \textbar\ Author's Claim}
  \examplemeta{Indikator: inference from evidence}{Kesulitan: tinggi}{Target: 100 detik}
  \begin{block}{Passage}
    \fontsize{9.8}{12.1}\selectfont Students often treat a practice-test score as a fixed measure of ability. Yet the score combines several factors: conceptual knowledge, pacing, fatigue, and even the order in which questions are attempted. A useful review therefore goes beyond counting correct answers. It identifies which factor caused each error and selects a response that matches that cause.
  \end{block}
  {\bfseries Which statement best reflects the author's argument?}\par
  \vspace{0.06cm}
  {\fontsize{9.0}{10.9}\selectfont
  A. Practice tests accurately measure permanent academic ability.\par\vspace{0.035cm}
  B. Students should always answer questions in the printed order.\par\vspace{0.035cm}
  C. Reviewing errors by cause produces more useful follow-up than using the total score alone.\par\vspace{0.035cm}
  D. Fatigue is the main reason students answer incorrectly.\par\vspace{0.035cm}
  E. Conceptual study should be replaced by repeated tests.}
\end{frame}

\begin{frame}{Bedah Contoh 6 -- LBE}
  \kicker{Jawaban dan Proses Berpikir}
  \answer{C. Review by cause is more useful than the total score alone.}
  \vspace{0.20cm}
  \begin{block}{Why?}
    The passage lists several sources of error and concludes that follow-up should match the specific cause. Option C restates this central argument without adding an unsupported absolute claim.
  \end{block}
  {\fontsize{10.6}{13.2}\selectfont
  \mistake{choosing an option that repeats one familiar word from the passage but changes the strength of the claim.}\par\vspace{0.16cm}
  \quick{summarize each sentence in a few words, then choose the option that connects the entire paragraph.}}
\end{frame}

\begin{frame}[shrink=12]{Contoh 7 -- Penalaran Matematika}
  \kicker{Soal \textbar\ Kecukupan Data}
  \examplemeta{Indikator: merumuskan model}{Kesulitan: tinggi}{Target: 110 detik}
  \begin{block}{Kasus}
    Satu sesi literasi berlangsung 35 menit dan satu sesi kuantitatif berlangsung 50 menit. Berapa banyak sesi literasi dalam suatu pekan?
  \end{block}
  {\fontsize{10.0}{12.3}\selectfont
  (1) Jumlah seluruh sesi pada pekan itu adalah 17.\par
  (2) Jumlah waktu seluruh sesi pada pekan itu adalah 700 menit.}\par
  \vspace{0.10cm}
  {\fontsize{9.15}{11.2}\selectfont
  A. Pernyataan (1) saja cukup, sedangkan (2) saja tidak cukup.\par\vspace{0.035cm}
  B. Pernyataan (2) saja cukup, sedangkan (1) saja tidak cukup.\par\vspace{0.035cm}
  C. Masing-masing pernyataan cukup.\par\vspace{0.035cm}
  D. Kedua pernyataan bersama-sama cukup, tetapi masing-masing tidak cukup.\par\vspace{0.035cm}
  E. Kedua pernyataan bersama-sama tidak cukup.}
\end{frame}

\begin{frame}[shrink=6]{Bedah Contoh 7 -- Penalaran Matematika}
  \kicker{Jawaban dan Proses Berpikir}
  \answer{D. Kedua pernyataan bersama-sama cukup.}
  \vspace{0.16cm}
  {\fontsize{10.5}{13.0}\selectfont Misalkan $l$ adalah sesi literasi dan $q$ sesi kuantitatif.
  \begin{columns}[T,onlytextwidth]
    \begin{column}{0.47\textwidth}
      \begin{block}{Pernyataan (1)}
        $l+q=17$ memiliki banyak pasangan; belum cukup.
      \end{block}
    \end{column}
    \begin{column}{0.47\textwidth}
      \begin{block}{Pernyataan (2)}
        $35l+50q=700$ memiliki lebih dari satu pasangan; belum cukup.
      \end{block}
    \end{column}
  \end{columns}
  Bersama-sama, $l+q=17$ dan $35l+50q=700$ menghasilkan $l=10$ dan $q=7$.\par\vspace{0.11cm}
  \mistake{langsung menyelesaikan dua persamaan tanpa memeriksa kecukupan setiap pernyataan secara terpisah.}\par\vspace{0.11cm}
  \quick{uji (1), uji (2), baru uji gabungan; yang dinilai adalah kecukupan, bukan sekadar hasil akhir.}}
\end{frame}

\begin{frame}{Kesalahan literasi mengarahkan latihan membaca}
  \kicker{Dari Jawaban Menuju Intervensi -- Bagian 1}
  \begin{columns}[T,onlytextwidth]
    \begin{column}{0.30\textwidth}
      \begin{block}{Stimulus keliru}
        Informasi belum terbaca. Latih anotasi dan parafrasa.
      \end{block}
    \end{column}
    \begin{column}{0.30\textwidth}
      \begin{block}{Opsi mutlak}
        Kekuatan klaim terlewat. Tandai kata \emph{selalu}, \emph{pasti}, dan \emph{hanya}.
      \end{block}
    \end{column}
    \begin{column}{0.30\textwidth}
      \begin{block}{Makna tersirat}
        Bukti belum terhubung. Latih simbol, rujukan, dan inferensi.
      \end{block}
    \end{column}
  \end{columns}
  \vspace{0.18cm}
  \claim{Opsi yang salah tetap bernilai: ia menunjukkan bagian proses yang perlu diperbaiki.}
\end{frame}

\begin{frame}{Kesalahan numerasi mengarahkan latihan strategi}
  \kicker{Dari Jawaban Menuju Intervensi -- Bagian 2}
  \begin{columns}[T,onlytextwidth]
    \begin{column}{0.30\textwidth}
      \begin{block}{Benar, terlalu lama}
        Strategi belum efisien. Bandingkan dua cara penyelesaian.
      \end{block}
    \end{column}
    \begin{column}{0.30\textwidth}
      \begin{block}{Model keliru}
        Konteks belum diterjemahkan. Latih variabel dan batasan.
      \end{block}
    \end{column}
    \begin{column}{0.30\textwidth}
      \begin{block}{Salah berulang}
        Konsep belum tuntas. Kembali ke modul dan remedial topikal.
      \end{block}
    \end{column}
  \end{columns}
  \vspace{0.18cm}
  \claim{Tindak lanjut dipilih dari penyebab kesalahan, bukan dari nilai total saja.}
\end{frame}

\begin{frame}{Orang tua dapat mendukung proses tanpa mengajar isi soal}
  \kicker{Peran di Rumah}
  \begin{columns}[T,onlytextwidth]
    \begin{column}{0.46\textwidth}
      {\color{StudentBlue}\bfseries Pertanyaan yang membantu}\par
      \begin{itemize}
        \item Bagian mana yang membuatmu ragu?
        \item Apa bukti untuk jawabanmu?
        \item Kesalahan ini konsep, strategi, atau ketelitian?
        \item Apa rencana latihan berikutnya?
      \end{itemize}
    \end{column}
    \begin{column}{0.48\textwidth}
      {\color{StudentBlue}\bfseries Dukungan yang dibutuhkan}\par
      \begin{itemize}
        \item Waktu belajar yang konsisten.
        \item Ruang untuk membahas kesalahan.
        \item Apresiasi pada perbaikan proses.
        \item Komunikasi dengan pembimbing berbasis data.
      \end{itemize}
    \end{column}
  \end{columns}
  \vspace{0.15cm}
  \begin{beamercolorbox}[rounded=true,sep=0.18cm,wd=\textwidth]{block body}
    \centering\bfseries Tujuannya bukan mengejar jawaban cepat di rumah, melainkan membangun kebiasaan refleksi.
  \end{beamercolorbox}
\end{frame}

\begin{frame}{Kesimpulan contoh soal UTBK}
  \kicker{Pesan Utama Bab 7}
  \begin{columns}[T,onlytextwidth]
    \begin{column}{0.30\textwidth}
      \begin{block}{Beragam}
        Stimulus dan bentuk tugas berbeda pada setiap subtes.
      \end{block}
    \end{column}
    \begin{column}{0.30\textwidth}
      \begin{block}{Bernalar}
        Jawaban harus didukung bukti, hubungan, atau model yang sah.
      \end{block}
    \end{column}
    \begin{column}{0.30\textwidth}
      \begin{block}{Bertumbuh}
        Kesalahan dianalisis untuk memilih latihan yang tepat.
      \end{block}
    \end{column}
  \end{columns}
  \vspace{0.26cm}
  \centering{\color{StudentNavy}\bfseries\fontsize{14}{17}\selectfont Bab 7 adalah contoh cara berpikir, bukan bank soal penuh.}
\end{frame}

\begin{frame}{Sumber resmi dan dasar penyusunan}
  \kicker{Rujukan Bab 7 -- Bagian 1}
  {\color{StudentBlue}\bfseries\fontsize{14.5}{17.5}\selectfont 1. Panitia SNPMB}\par
  \vspace{0.05cm}
  {\fontsize{11}{13.5}\selectfont Framework UTBK--SNBT Tahun 2026.}\par
  {\fontsize{10.5}{13}\selectfont\href{https://snpmb.id/fr/}{https://snpmb.id/fr/}}\par
  {\color{StudentGray}\fontsize{10}{12.4}\selectfont Acuan kemampuan, proses, konteks, jumlah soal, dan durasi tujuh subtes.}\par
  \vspace{0.34cm}
  {\color{StudentBlue}\bfseries\fontsize{14.5}{17.5}\selectfont 2. Panitia SNPMB}\par
  \vspace{0.05cm}
  {\fontsize{11}{13.5}\selectfont Informasi Umum UTBK--SNBT 2026.}\par
  {\fontsize{10.5}{13}\selectfont\href{https://snpmb.id/utbk-snbt/informasi-umum}{snpmb.id/utbk-snbt/informasi-umum}}\par
  {\color{StudentGray}\fontsize{10}{12.4}\selectfont Digunakan untuk menjaga batas antara informasi resmi 2026 dan persiapan 2027.}
\end{frame}

\begin{frame}{Status contoh dan penggunaan}
  \kicker{Rujukan Bab 7 -- Bagian 2}
  {\color{StudentBlue}\bfseries\fontsize{14.5}{17.5}\selectfont 3. Tim Akademik StudentPro}\par
  \vspace{0.05cm}
  {\fontsize{10.8}{13.3}\selectfont Contoh soal, opsi, kunci, pembahasan, analisis kesalahan, dan strategi cepat pada bab ini disusun secara orisinal untuk kebutuhan sosialisasi dan pembelajaran.}\par
  \vspace{0.28cm}
  {\color{StudentBlue}\bfseries\fontsize{14.5}{17.5}\selectfont 4. Catatan batas penggunaan}\par
  \vspace{0.05cm}
  {\fontsize{10.8}{13.3}\selectfont Contoh tidak mereproduksi soal resmi, tidak memprediksi soal yang akan keluar, dan tidak menetapkan format final UTBK 2027.}\par
  \vfill
  \begin{beamercolorbox}[rounded=true,sep=0.17cm,wd=\textwidth]{block body}
    \centering\bfseries Bank soal, latihan bertingkat, dan tryout dikelola terpisah dalam program StudentPro.
  \end{beamercolorbox}
\end{frame}
